\section{Gluing a left state through its boundary contraction}

This section records the boundary contractions and the gluing lemma of Landau,
Vazirani and Vidick~\cite[lines 409--447]{Landau2013PolynomialTime}. Fix a cut
of $(\mathcal A\otimes\mathcal S)\otimes\mathcal R$, where $\mathcal S$ is one
site, $\mathcal A$ the part strictly to its left and $\mathcal R$ the part to
its right, and let $W:\mathbb C^B\to\mathcal R$ be an isometry. In the source,
$W=U_v$ maps $\lvert j\rangle$ to the $j$-th right Schmidt vector of a vector
$v$ of Schmidt rank $B$.

\begin{definition}[Left state and boundary contraction]
    \label{def:boundary_contraction}
    \lean{Matrix.partialTraceFirst, Matrix.conjTranspose_mulVec_gluing}
    \leanok
    For a matrix $X$ on $(\mathcal A\otimes\mathcal S)\otimes\mathcal C$, write
    $\operatorname{tr}_{\mathcal A}X$ for its partial trace over
    $\mathcal A$, a matrix on $\mathcal S\otimes\mathcal C$. If
    $v=(1\otimes W)x$, then $x=(1\otimes W)^\dagger v$ is the left state of
    $v$, and $\operatorname{tr}_{\mathcal A}\lvert x\rangle\langle x\rvert$ is
    its boundary contraction.
\end{definition}

\begin{theorem}[Gluing]
    \label{thm:boundary_gluing}
    \lean{Matrix.partialTraceRight_gluing, Matrix.partialTraceFirst_gluing,
        Matrix.rectangularTraceNorm_partialTraceFirst_gluing_sub,
        Matrix.norm_trace_mul_le_l2_opNorm_mul_rectangularTraceNorm,
        Matrix.re_trace_gluing_le}
    \leanok
    \uses{def:boundary_contraction}
    Let $W$ be an isometry, $\sigma$ a matrix on
    $(\mathcal A\otimes\mathcal S)\otimes\mathbb C^B$,
    $\sigma'=(1\otimes W)\sigma(1\otimes W)^\dagger$ and $v=(1\otimes W)x$.
    Then
    \begin{enumerate}
        \item $\operatorname{tr}_{\mathcal R}\sigma'
                  =\operatorname{tr}_{\mathbb C^B}\sigma$;
        \item $\lVert\operatorname{tr}_{\mathcal A}\sigma'
                  -\operatorname{tr}_{\mathcal A}\lvert v\rangle\langle v\rvert\rVert_1
                  =\lVert\operatorname{tr}_{\mathcal A}\sigma
                  -\operatorname{tr}_{\mathcal A}\lvert x\rangle\langle x\rvert\rVert_1$;
        \item for $H=H_L\otimes1+1\otimes K$, with $H_L$ acting on
              $\mathcal A\otimes\mathcal S$ and $K$ on
              $\mathcal S\otimes\mathcal R$,
              \begin{align*}
                  \operatorname{Re}\operatorname{tr}(\sigma'H)
                   & \le\operatorname{Re}\operatorname{tr}(\sigma(H_L\otimes1))
                  +\operatorname{Re}\langle v,(1\otimes K)v\rangle
                  +\lVert K\rVert\,\lVert\operatorname{tr}_{\mathcal A}\sigma
                  -\operatorname{tr}_{\mathcal A}\lvert x\rangle\langle x\rvert\rVert_1.
              \end{align*}
    \end{enumerate}
    This is Lemma~l:gluing of~\cite[lines 421--447]{Landau2013PolynomialTime},
    for every isometry $W$ whose range contains the right Schmidt vectors of
    $v$. The constant $n$ of the printed third part is replaced by $\lVert K\rVert$,
    which is at most $n$ for the chain Hamiltonian of the source, where
    $K=H_i+H_R$.
\end{theorem}

\begin{proof}
    \leanok
    The first part holds because $W^\dagger W=1$. For any matrix $X$,
    $\operatorname{tr}_{\mathcal A}((1\otimes W)X(1\otimes W)^\dagger)
        =(1\otimes W)(\operatorname{tr}_{\mathcal A}X)(1\otimes W)^\dagger$, and
    $\lvert v\rangle\langle v\rvert
        =(1\otimes W)\lvert x\rangle\langle x\rvert(1\otimes W)^\dagger$; the
    second part follows since conjugation by an isometry preserves the trace
    norm. For the third, $(1\otimes W)^\dagger(H_L\otimes1)(1\otimes W)
        =H_L\otimes1$, and
    $\operatorname{tr}(\sigma'(1\otimes K))
        -\langle v,(1\otimes K)v\rangle
        =\operatorname{tr}\bigl((\operatorname{tr}_{\mathcal A}\sigma'
        -\operatorname{tr}_{\mathcal A}\lvert v\rangle\langle v\rvert)K\bigr)$,
    whose modulus is at most $\lVert K\rVert$ times the trace norm of the
    difference, which is the quantity of the second part.
\end{proof}

\begin{theorem}[A dominant low-energy eigenvector]
    \label{thm:dominant_low_energy_eigenvector}
    \lean{Matrix.IsHermitian.re_trace_mul_eq_sum_eigenvalues_mul,
        Matrix.PosSemidef.exists_eigenvector_of_re_trace_mul_le}
    \leanok
    Let $H-E_0\ge\Delta(1-\lvert\Gamma\rangle\langle\Gamma\rvert)$ with
    $\Delta>0$ and $\lVert\Gamma\rVert=1$, and let $\sigma$ be a density matrix
    with $\operatorname{Re}\operatorname{tr}(\sigma H)\le E_0+e$, where
    $4e<\Delta$. Then $\sigma$ has a unit eigenvector $v$ with eigenvalue larger
    than $\tfrac12$, energy
    $\operatorname{Re}\langle v,Hv\rangle\le E_0+2e$ and overlap
    $\lvert\langle\Gamma,v\rangle\rvert\ge1-2e/\Delta$.
\end{theorem}

\begin{proof}
    \leanok
    Write $\sigma=\sum_k\lambda_k\lvert v_k\rangle\langle v_k\rvert$, so that
    $\sum_k\lambda_k(\langle v_k,Hv_k\rangle-E_0)\le e$ with every term
    nonnegative. The eigenvectors of energy above $E_0+2e$ therefore carry total
    weight less than $\tfrac12$. Each eigenvector of energy at most $E_0+2e$ has
    $\lvert\langle\Gamma,v_k\rangle\rvert^2\ge1-2e/\Delta>\tfrac12$, and since
    $\sum_k\lvert\langle\Gamma,v_k\rangle\rvert^2=1$ there is at most one of
    them; its eigenvalue exceeds $\tfrac12$, and the overlap is at least its
    square.
\end{proof}

\begin{theorem}[The leading eigenvector of a near-optimal left state]
    \label{thm:leading_eigenvector_gluing}
    \lean{Matrix.trace_gluing_mul_kronecker_one,
        Matrix.exists_leading_eigenvector_of_gluing}
    \leanok
    \uses{thm:boundary_gluing, thm:dominant_low_energy_eigenvector}
    In the setting of Theorem~\ref{thm:boundary_gluing}, let
    $H-E_0\ge\Delta(1-\lvert\Gamma\rangle\langle\Gamma\rvert)$ with
    $\Delta>0$ and $\lVert\Gamma\rVert=1$, and let $v=(1\otimes W)x$ have
    $\operatorname{Re}\langle v,Hv\rangle\le E_0+e$. Let $\sigma$ be a density
    matrix on $(\mathcal A\otimes\mathcal S)\otimes\mathbb C^B$ with
    $\operatorname{Re}\operatorname{tr}(\sigma(H_L\otimes1))\le
        \operatorname{Re}\operatorname{tr}(\lvert x\rangle\langle x\rvert(H_L\otimes1))$
    and
    $\lVert\operatorname{tr}_{\mathcal A}\sigma
        -\operatorname{tr}_{\mathcal A}\lvert x\rangle\langle x\rvert\rVert_1\le\tau$.
    If $4(e+\lVert K\rVert\tau)<\Delta$, then $\sigma$ has a unit eigenvector
    $u$ with eigenvalue larger than $\tfrac12$ such that $u'=(1\otimes W)u$ has
    energy at most $E_0+2(e+\lVert K\rVert\tau)$ and
    $\lvert\langle\Gamma,u'\rangle\rvert\ge1-2(e+\lVert K\rVert\tau)/\Delta$.

    This is Lemma~l:leading of~\cite[lines 493--505]{Landau2013PolynomialTime},
    for the solution $\sigma$ of the size-trimming program: its left state $x$ of
    $v$ is feasible, so $\sigma$ has left energy at most that of $x$, and
    $\tau=c_\varepsilon/n$. With $\lVert K\rVert\le n$, $e$ is the energy
    excess of the trimmed state of Lemma~l:within, whose printed value
    $6\sqrt{c_\varepsilon}$ is not established by the printed argument; the
    corrected excess and a choice of $c_\varepsilon$ with
    $2(e+c_\varepsilon)<\varepsilon/12$ are given
    in~\cite{gap:lvv13_trimming_energy}.
\end{theorem}

\begin{proof}
    \leanok
    By Theorem~\ref{thm:boundary_gluing} and the left-energy bound,
    $\sigma'=(1\otimes W)\sigma(1\otimes W)^\dagger$ has energy at most
    $E_0+e+\lVert K\rVert\tau$, and it is a density matrix. Theorem~\ref{thm:dominant_low_energy_eigenvector}
    gives an eigenvector $v'$ of $\sigma'$ with eigenvalue $\mu>\tfrac12$. Since
    $\mu\ne0$, $v'=\mu^{-1}\sigma'v'$ lies in the range of $1\otimes W$, so
    $u=(1\otimes W)^\dagger v'$ satisfies $(1\otimes W)u=v'$ and
    $\sigma u=\mu u$.
\end{proof}
